\section{Conclusions}

\begin{frame}{Conclusions}
	In this course, we saw:
	\begin{itemize}
		\item That Principal Component Analysis is an \textbf{optimal} solution to the problem of \textbf{linearly} combining input features in order to maximize variance.
		\pause
		\item That Principal Component Analysis allows us to \textbf{redistribute} the variance of the features \textbf{without loss}, while rendering the linear combinations completely uncorrelated.
		\pause
		\item That selecting the top $k$ largest principal components allows us to \textbf{retain} a large portion of the \textbf{variance} while \textbf{reducing} the \textbf{dimensionality} of the problem.
		\pause
		\item How a machine learning algorithm uses images as input data.
		\pause
		\item How Principal Component Analysis successfully simplified images from a handwritten digit dataset by reducing the data size by a factor of 5 while preserving 95\% of the variance.
	\end{itemize}
\end{frame}